\section{Algebra}
\label{sec:algebra}

The \texttt{primepie.algebra} module implements basic algebraic structures
used throughout the library, with performance-critical paths implemented in C++
and exposed to Python via \texttt{pybind11}.

\subsection{Algebraic Hierarchy}

At this stage the library provides two families of structures:

\begin{itemize}
  \item \textbf{Groups}: an abstract interface with a binary operation,
        identity, and inverse. (In our Python API, the additive group
        of $\mathbb{Z}/n\mathbb{Z}$ is exposed.)
  \item \textbf{Rings}: extend groups by adding a multiplication together with
        \texttt{one()} and \texttt{mul()}. We currently expose the ring of
        integers $\mathbb{Z}$ and the modular rings $\mathbb{Z}/n\mathbb{Z}$.
\end{itemize}

Both C++ and Python sides offer consistent scalar and batched operations.

\subsection{Exponentiation in Groups and Rings}

Exponentiation is implemented generically via
\texttt{GroupUtils::power} (additive/group law) and
\texttt{RingUtils::power} (multiplicative law). Both use classical
binary exponentiation, computing $a^n$ in $\mathcal{O}(\log n)$ steps by
repeated squaring:
\[
a^n =
\begin{cases}
1, & n = 0,\\
(a^{n/2})^2, & n \text{ even},\\
a \cdot (a^{(n-1)/2})^2, & n \text{ odd}.
\end{cases}
\]
In the group context, this means repeated application of the group operation
(addition in additive groups). In the ring context, it means repeated
multiplication (requiring $n \ge 0$ in general rings).

\subsection{The Ring of Integers \texorpdfstring{$\mathbb{Z}$}{Z}}

The class \texttt{Integers} models $(\mathbb{Z}, +, \cdot)$ and exposes
canonical operations:
\[
\begin{aligned}
a + b &= \texttt{add}(a,b), \qquad
a \cdot b = \texttt{mul}(a,b), \\
-a &= \texttt{neg}(a), \qquad
0 = \texttt{zero}(), \; 1 = \texttt{one}().
\end{aligned}
\]
These are implemented with native integer arithmetic. The unified \texttt{Ring}
interface also underlies modular rings~$\mathbb{Z}/n\mathbb{Z}$.

\subsubsection*{Greatest Common Divisor}

\texttt{Integers::gcd} implements the Euclidean algorithm, returning a
nonnegative result and obeying
$\gcd(0,0)=0$, $\gcd(a,0)=|a|$, and $\gcd(0,b)=|b|$.
The core loop is the standard reduction
\[
\texttt{while (b != 0)} \{ \; r = a \bmod b;\; a = b;\; b = r; \;\}.
\]
For batches, the same computation is applied elementwise (optionally in
parallel when compiled with OpenMP).

\subsubsection*{Deterministic Primality Testing}

\texttt{Integers::is\_prime} is based on the Miller–Rabin test
\cite{miller1976rabin}, made deterministic for all 64-bit integers by using a
fixed set of bases. Small factors are removed by trial division over
$\{2,3,5,7,11,13,17,19,23,29,31,37\}$ and then Miller–Rabin is run with that
same set. For $n < 2^{64}$ this set of bases is sufficient to guarantee
determinism \cite{jaeschke1993miller, sorenson2015strong}.

